%% notes-tex/eqtl-data-model/sections/3-ontology.tex
%% [[torchcell.datamodels.eqtl-data-model]]
%% https://github.com/Mjvolk3/torchcell/tree/main/notes-tex/eqtl-data-model/sections/3-ontology.tex
%% SOURCE: hand-authored. Every class name and the systematic-gene-name pattern are read
%% from torchcell/datamodels/schema.py, not recalled.

\section{How it fits the ontology\secstatus{todo}}\label{sec:ontology}

\subsection{What already exists\secstatus{todo}}

The primary record is a genotype and an environment mapped to a phenotype, which is exactly
\file{Experiment}. Nothing conceptual is missing, and four pieces are already built.

\begin{itemize}[leftmargin=1.2em,itemsep=2pt,topsep=2pt]
  \item \file{RNASeqExpressionPhenotype} is the bulk readout, and it was written for this
    situation. Its docstring describes a survey ``where each isolate's transcriptome is
    measured on its OWN genome,'' with ``no common reference to ratio against,'' storing
    absolute TPM. Albert 2018 and Smith 2008 fit unchanged.
  \item \file{PseudobulkExpressionPhenotype} is the single-cell readout, built for
    Nadal-Ribelles 2025. It carries \file{dispersion} and \file{n_cells} per genotype,
    which is the right shape for Boocock's noise eQTLs: heterogeneity survives as a
    per-genotype scalar rather than being averaged away.
  \item \file{SequenceVariantPerturbation} states the intent outright. Together with
    \file{CopyNumberVariantPerturbation} it ``lets a NATURAL ISOLATE be modeled as a
    perturbation set off S288C ($\sim$4,500 sequence variants per isolate).''
  \item The off-graph pointer pattern on that class, \file{sequence_source},
    \file{sequence_uri} and \file{sequence_sha256}, means variant sequences are never
    inlined.
\end{itemize}

\subsection{Where the obvious mapping fails\secstatus{todo}}\label{sec:fails}

Reusing \file{SequenceVariantPerturbation} for a segregant is the obvious move and it is
wrong. \textbf{A natural isolate and a segregant are not the same kind of object.} Caudal's
isolates are each individually sequenced, so the variant set is observed and every allele
dereferences to a real sequence with a hash. A segregant is never sequenced to that
standard; the \emph{parents} are, and the segregant is imputed as a mosaic of them. That
class promises a dereferenceable, hash-verified sequence, and for a segregant no such
artifact exists. Using it would encode an inference as an observation, which is the precise
failure the verification levels exist to prevent.

\subsection{Four gaps, in order of how much they matter\secstatus{todo}}

\notesec{Intergenic variants have no representation} \file{GenePerturbation} requires
\file{systematic_gene_name} and validates it against

\begin{quote}\small
\verb+^(Y[A-P][LR]\d{3}[WC](-[A-Z])?|Q\d{4}|YNC[A-Q]\d{4}[WC])$+
\end{quote}

so every perturbation must attach to a real gene. A promoter SNP 300 bp upstream of an ORF
has no gene to attach to, and assigning it to the downstream gene is an interpretation,
wrong whenever the variant acts on a divergent neighbor. This is the sharpest gap and it is
not incidental: the promoter is where the strongest and most interpretable cis signal lives.
It is also the same gap that blocks every massively parallel reporter assay, which perturbs
regulatory sequence exclusively. A gene-anchored patch does exist: the genome machinery
already carries a flanking-window convention for gene sequence (1,000 bp upstream, 300 bp
downstream), which covers most promoter variants, and the anchor need not even be a
natural gene, since a synthetic locus name tied to a sequence hash could stand in for an
edited or renamed region. Either way the window assignment is an interpretation and would
have to be recorded as one.

\notesec{Cardinality} An SGA record carries at most three perturbations. A segregant carries
$\sim\!10^4$. For $n \sim 10^3$ segregants that is $\sim\!10^7$ perturbation objects for one
dataset, and \file{Genotype} sorts and set-compares its perturbation list on every equality
check while the L1 key hashes the genotype signature.

\notesec{Reference asymmetry} Reference here means the genome that variants are expressed
against, not the baseline a phenotype is ratioed to; a record has both, and this gap
concerns only the first. A segregant's natural genomic reference is one of its two
parents, but the convention expresses every genotype off S288C. BY is near-isogenic to
S288C, so BY-derived blocks are nearly empty perturbation sets while RM-derived blocks
carry the full variant load. Keeping S288C is still right, since it is what makes these
records comparable to everything else, but the lopsidedness reflects the reference choice
rather than the biology and should be documented rather than discovered. The phenotype
side is untouched by this: expression is stored absolute, with no wild-type ratio in
these designs at all.

\notesec{Linkage is destroyed by a set} Variants arrive in blocks and co-segregate.
\file{Genotype} does hold its perturbations in a stable order, but that order is a
canonical sort by systematic gene name and perturbation type, imposed so that equality
and hashing are well defined; it is set semantics with deterministic iteration, and
lexical gene-name order is not chromosomal order. Nothing in the representation says two
variants sit on the same block and travel together, which is exactly the structure that
makes mapping possible and that limits attribution.

\subsection{The representation that follows\secstatus{todo}}\label{sec:proposal}

Store what was actually determined, a \term{haplotype mosaic} rather than a variant set:

$$
G_i \;=\; \big\{ (\text{chr},\, \text{start},\, \text{end},\, \text{parent},\, p) \big\},
\qquad |G_i| \sim 10^1\text{--}10^2
$$

against two hash-pinned parent assemblies, with a per-interval posterior $p$ and the marker
spacing recorded as the resolution. This is better on four counts at once.

\begin{enumerate}[leftmargin=1.4em,itemsep=1pt,topsep=2pt]
  \item \textbf{Honest.} It stores the inference as an inference, with its uncertainty,
    rather than promising sequence that was never read.
  \item \textbf{Compact.} Tens of intervals per strain instead of $\sim\!10^4$ variants,
    which dissolves the cardinality problem rather than solving it.
  \item \textbf{Lossless.} The full variant set is recoverable exactly from the two parent
    assemblies plus the intervals, so it becomes a derived view and nothing is given up.
  \item \textbf{Linkage-preserving.} The interval is the unit, so co-segregation is
    explicit.
\end{enumerate}

The cost is a genuinely new perturbation type, keyed on a genomic interval and a parent
rather than on a gene, carrying a probability. That cost is accepted: the type is judged
necessary, not merely convenient. It is also the same non-gene-keyed requirement the
promoter-variant gap already imposes, so the two should be resolved together.

\subsection{What to do\secstatus{todo}}\label{sec:todo}

\begin{enumerate}[leftmargin=1.4em,itemsep=2pt,topsep=2pt]
  \item \textbf{Ingest the two matrices, not the QTL table.} One
    \file{RNASeqExpressionExperiment} per segregant per condition, or
    \file{PseudobulkExpressionExperiment} for the single-cell studies. The QTL table is an
    estimate conditional on method and threshold, it has no source value to check against
    at any verification level, and storing it as a \file{Phenotype} would be ingesting
    somebody's analysis as data.
  \item \textbf{Decide the intergenic question deliberately}, because it gates both this
    class and the regulatory-DNA class. With the mosaic type accepted, an interval-keyed
    perturbation leaf exists regardless, so the live options are: attach promoter
    variants to that same interval-keyed leaf; or keep gene-keying and assign each
    variant to a gene via the existing flanking-window convention (1,000 bp upstream,
    300 bp downstream), with the attribution recorded as an interpretation, e.g. a
    \file{ProvenanceGap}; or restrict to coding variants. The third is the honest
    short-term option and the worst long-term one.
  \item \textbf{Treat the parent assemblies as an ingestion dependency.} A segregant record
    means nothing without them, so they are part of the artifact rather than context, and
    each needs a hash-pinned assembly rather than a variant call file.
  \item \textbf{Measure the cardinality cost} on a single Albert 2018 segregant before
    building anything.
\end{enumerate}

\subsection{What was built\secstatus{todo}}\label{sec:built}

The mosaic was implemented on 2026-09-12 for the Bloom 2019 segregant panels (eLife
8:e49212, doi 10.7554/eLife.49212; 16 crosses, 13,950 haploid segregants, 38 growth
conditions), in \file{torchcell/datasets/scerevisiae/bloom2019.py}, and it departs from
Sec.~\ref{sec:proposal} in one respect. The proposal called the cost ``a genuinely new
perturbation type''. A perturbation leaf cannot carry a mosaic: \file{Genotype} sorts its
perturbations on \file{systematic_gene_name} and every gene-keyed consumer walks that
list, so a leaf with no gene either fails the regex or has to carry a name it does not
have, and a \file{Genotype} with an empty list reads as wild-type S288C to every consumer
that reads only the list. The mosaic is therefore a \emph{sibling} of \file{Genotype}:
\file{SegregantGenotype} holds \file{HaplotypeBlock}s of (chromosome, start, end, parent,
posterior, n\_markers) between two \file{SegregantParent}s pinned to sha256-anchored
assemblies, and \file{SegregantGrowthExperiment} declares that class as its genotype.
The base \file{Experiment} rejects a segregant dict, so the union resolves the class
unambiguously and a gene-keyed consumer fails loudly rather than reading a segregant as
unperturbed. The new classes enter only the module-level unions and maps, so no served
dataset's schema contract moved and the admission gate reported the increment admissible.

Three of the four properties in Sec.~\ref{sec:proposal} held as stated on the built
store. Honest: the release is the R/qtl hard call, so every block carries posterior 1.0
and the call method is quoted from the authors' code. Lossless: every one of the 13,950
block sets re-expands to its released marker row, checked at L2. Linkage-preserving by
construction. Compact needs a correction. Cross A, the BY $\times$ RM cross the estimate
of Sec.~\ref{sec:proposal} was made for, has a median of 83 blocks per segregant, inside
the $10^1$--$10^2$ range; the other 15 crosses run from 90 to 143 at the median, and the
means are pulled to 100--240 by tails of segregants with hundreds to 1,847 blocks. A
segregant is still a few kilobytes rather than $\sim10^4$ variants, so the cardinality
argument stands, but the block count is not a property of meiosis alone. It is a
hypothesis, untested here, that the tail is short runs from genotyping noise in the hard
calls that a posterior threshold would collapse; the release carries no posteriors, so the
blocks record the calls as released.

Two items of Sec.~\ref{sec:todo} closed with the build and two stay open. The parent
assemblies are an ingestion dependency: each parent is pinned to its member of the 1011
assemblies tarball, and a parent absent from the member index raises. The cardinality
cost was measured on the whole panel rather than on one Albert 2018 segregant. The
intergenic question is not resolved but is no longer coupled to this class, since a
mosaic is not a \file{GenePerturbation} and needs no gene; a cis-eQTL still has no gene
to attach to, so per-variant attribution and the regulatory-DNA class remain one open
decision. Ingesting the two matrices for Albert 2018 and Boocock 2025 is the natural next
increment on the same genotype class.
